\section{结论与下一步验证}
本文给出一条以题目附件为主、在证据断点处停止数值传递的资源配置分析链。第一问的稳健评分揭示领域间相对差异和 c4 指标冲突；二次 ILR 对 1M 配比 Loss 有可检验预测力，但实测优良配方质心仍是待训练验证的候选。第二问的 B1 幂律和 B6 质量修正分别在各自数据内成立；选中结构表明半合成质量水平可改变参数与数据项的条件 Loss，但不能证明第一问质量评分具有相同效应。第三问在固定候选配比和题设成本下得到预算路径；低预算受 $D$ 下界影响，高预算出现支持域饱和。第四问可描述严格开放样本的观测能力和条件算力路径，却不能可靠估计 12/24 个月能力数值。

最直接的后续实验是：在同一训练与评测协议下，对多个 $N,D$、至少两种 17 域配比、多个可复测的质量水平作成对训练；同时记录真实质量处理成本与上下文长度。这样才能标定 $Q_A\leftrightarrow Q_B$、比较配比与质量的交互结构，并把前三问的 Loss 决策桥接至第四问能力指标。没有这些共同实验前，本文不把条件最优和历史关联包装为现实的因果或长期预测结论。
